%%%%%%%%%%%%%%%%%%%%%%%%%%%%%%%%%%%%%%%%%%%%%%%%%%%%%
%			GENERÁTOR v. 18.09 				%
%%%%%%%%%%%%%%%%%%%%%%%%%%%%%%%%%%%%%%%%%%%%%%%%%%%%%
% Tento soubor slouží k jednotlivé zkusné kompilaci 
% písniček či k kompilaci samostatných pdf souborů 
% pomocí bashového skriptu kompilace.sh ve stejné 
% složce. Z důvodu hromadné kompilace musí poslední 
% dva řádky tohoto souboru zůstat nezměněny.
%%%%%%%%%%%%%%%%%%%%%%%%%%%%%%%%%%%%%%%%%%%%%%%%%%%%%
%			Jak kompilovat jednotlivé písně?        %
%%%%%%%%%%%%%%%%%%%%%%%%%%%%%%%%%%%%%%%%%%%%%%%%%%%%%
%	1. Do míst PÍSNIČKA upravte následnující řádek, 
%	   aby vypadal takto: 
%	   \input{../songy/[JMÉNO SOUBORU KOMPILOVNÉ PÍSNIČKY].tex}
%	2. Soubor písničky musí být přeložitelný a musí 
%	   se nacházet ve složce ../songy.
%%%%%%%%%%%%%%%%%%%%%%%%%%%%%%%%%%%%%%%%%%%%%%%%%%%%%
%			Jak psát soubory songů?                 %
%%%%%%%%%%%%%%%%%%%%%%%%%%%%%%%%%%%%%%%%%%%%%%%%%%%%%
%	1. Více návodu je k tomuto napsáno v souboru 
%      ../songy/00Songtemplate. 
%%%%%%%%%%%%%%%%%%%%%%%%%%%%%%%%%%%%%%%%%%%%%%%%%%%%%
%			Jak kompilovat celý zpěvník?			%
%%%%%%%%%%%%%%%%%%%%%%%%%%%%%%%%%%%%%%%%%%%%%%%%%%%%%
%	1. Více návodu je k tomuto napsáno v souboru
%	   ../Cely_zpevnik/zpevnik.tex.
%%%%%%%%%%%%%%%%%%%%%%%%%%%%%%%%%%%%%%%%%%%%%%%%%%%%%

% Hlavička START
\documentclass[openany,12pt]{memoir}
\usepackage[utf8]{inputenc} 
\usepackage[czech]{babel}
\usepackage[T1]{fontenc}
\usepackage[top=1.5cm, bottom=2cm, left=2cm, right=2cm]{geometry}  % --> NASTAVENÍ OKRAJŮ
\usepackage{fancyhdr}
\usepackage{graphicx}
\usepackage{xwatermark}
\usepackage{xcolor}
\usepackage{changepage}
\usepackage{pdfpages}
\usepackage{lettrine}
\usepackage{indentfirst}  %Důležité pro formátování

%%%%%%%%%%%%%%%%%%%%%%%%%%%%%%%%%%%%%%
%  FONT                              %
%%%%%%%%%%%%%%%%%%%%%%%%%%%%%%%%%%%%%%
\usepackage{amssymb}
\usepackage{tgschola}



%%%%%% Package na zpěvník
\usepackage[full]{leadsheets}%http://mirrors.nic.cz/tex-archive/macros/latex/contrib/leadsheets/leadsheets_en.pdf   --> dokumentace	
\definesongtitletemplate{empty}{} 
\setchords{
format = \bfseries \sffamily,   %tučné akordy
minor = {mi},% 
input-notation = {german},%
output-notation = {german}%
}
\definesongtitletemplate{empty}{} 


\newlength{\drop}
% VODOZNAK
\newwatermark[pages=1-,color=red!50,angle=0,scale=2, xpos=0,ypos=0]{\includegraphics[width=5cm]{obr/pozadi2.jpg}} %--> dvojka na pozadí


%%%%%%%%%%%%%%%%%%%%%%%%%%%%%%%%%%%%%%%%%%%%%%%%%%
%		 Vlastní příkazy
\newcounter{Slokočet}   %Automatické číslování slok
\newcommand{\mezera}{
\phantom{.}

}   %Horizontální odsazení slok (poněkud blbě zadefinovaný, ale jinak se formát rozbije jako wtf prostě)
\newcommand{\stred}{5.2cm}   %%% Na zarovnání slok doprostřed, pozn. automatičtější zarovnávání na střed nejde
\newcommand{\carka}{,\:}
\newcommand{\m}[1]{\color{white}{#1}}  %Pro akordy
\newcommand{\ap}{'}	%Pro apostrof
\newcommand{\elipsa}{\kern\fontdimen3\font} %Příkaz pro lepší zacházení s výpustkami (=...); je to vpodstatě jen mezera mezi tečkama výpustky
\newcommand{\pindent}{17.62482 pt} %Správná velikost \parindentu u layoutu se dvěma minipageama
\newcommand{\predtitle}{\huge}
\newcommand{\mezisloupci}{\phantom{TT}} %Místo mezi dvěma sloupci na jedné stránce
\newcommand{\z}{\hspace*{\fill}\null}

%%% Možné velikosti písem 
\newcommand{\normalni}{\normalsize}
\newcommand{\velky}{\fontsize{14.4}{15}\selectfont}
\newcommand{\vetsi}{\fontsize{15}{16}\selectfont}
\newcommand{\nejvetsi}{\fontsize{16}{17}\selectfont}
\newcommand{\nejnejvetsi}{\fontsize{17}{19}\selectfont}

%%% Stará definice sloky spoléhající na indenty
%\newlength{\pismeno}
%\settowidth{\pismeno}{x} %Tohle není moc ideální velikost, ale funguje
%\newif\ifslokavelka
%\slokavelkafalse
%\newcommand{\sloka}{
%\ifnum \value{Slokočet}>8  %Pokud je sloka dvouciferná
%\mezera \noindent \addtocounter{Slokočet}{1} \hspace*{-\pismeno}\arabic{Slokočet}.
%\else %Pokud jen jednociferná
%\mezera \noindent \addtocounter{Slokočet}{1} \arabic{Slokočet}. 
%\fi
%} 	%sloka, která se automaticky čísluje

\newcommand{\distanc}{\:}  %Vzdálenost čísla sloky před slokou
\newlength{\delkaargumentu}
%%% Sloka s automatickým číslováním
\newcommand{\sloka}{%
\addtocounter{Slokočet}{1}% Zvýší se o 1 počet slok
\mezera%  Sloka se odsadí vertikálně
\settowidth{\delkaargumentu}{\arabic{Slokočet}.\distanc}% Zde se určí délka odsazení 
\hspace*{-\delkaargumentu}%
\arabic{Slokočet}.\distanc%
\ignorespaces% Aby nevznikaly zbytečné mezery
}


%%% Sloka s vlastním argumentem
\newcommand{\ssloka}[1]{%     
\settowidth{\delkaargumentu}{#1\distanc}
\mezera%
\hspace*{-\delkaargumentu}%
#1\distanc%
\ignorespaces%
}  

%%% Refrén
\newcommand{\refren}[1][0]{%  Nepovinný argument sděluje, kolikátý refrén toto je, bez argumentu se vytiskne pouze refrén
\ifnum #1>0 %Pokud nepovinný argument existuje
\mezera%
\settowidth{\delkaargumentu}{\textbf{R$_{\text{#1}}$:}\distanc}%
\hspace*{-\delkaargumentu}%
\textbf{R$_{\text{#1}}$:}\distanc%
\ignorespaces%
\else %Pokud nepovinný argument neexistuje
\mezera%
\settowidth{\delkaargumentu}{\textbf{R:}\distanc}%
\hspace*{-\delkaargumentu}%
\textbf{R:}\distanc%
\ignorespaces%
\fi
}

\newcommand{\predehra}{\ssloka{\textbf{Předehra:}}}


\addto\captionsczech{\renewcommand{\contentsname}{Seznam písní}}

%%%%%%%%%%%%%%%%%%%%%%%%%%%%%%%%%%%%
%    FORMÁTOVÁNÍ                   %
%%%%%%%%%%%%%%%%%%%%%%%%%%%%%%%%%%%%

%%% Vlevo zarovnaný text s blokem zarovnaným na střed
\usepackage{varwidth}% http://ctan.org/pkg/varwidth
\newenvironment{centerjustified}{%
  \begin{center} % so the minipage is centered
  \begin{varwidth}[t]{\textwidth}	
  \raggedright % so the minipage's text is left justified
  \setlength{\parindent}{\pindent}
}{%
  \end{varwidth}
  \end{center}
}


%Hlavnička END
\begin{document}
\sffamily %Sans-serif pro lepší čitelnost

\pagestyle{empty}
\newgeometry{top=1.5cm, bottom = 0cm, left = 2cm, right = 2cm}
	
\newpage

%Pro kompilaci je důležitý, aby ten \input a \end{document} zůstali na posledních
%dvou řádcích!!!!

\velky
\newpage
% PÍSNIČKA
\begin{song}{title=\predtitle\centering Slavíci z Madridu \\\large Waldemar Matuška \vspace*{-0.3cm}}  %% sem se napíše jméno songu a autor
\begin{centerjustified}

\refren[1]
	Na\elipsa\dots\,na na\elipsa\dots

\phantom{.}

\begin{minipage}{0.65\textwidth}
\sloka
	^{Ami\z}Nebe je modrý a ^{E\z}zlatý, ^{E\z}bílá sluneční ^{Ami\z}záře,\:\:

	^{Ami\z}horko a sváteční ^{E\z}šaty, ^{E\z}vřava a zpocený ^{Ami\z}tváře.

	^{Ami\z}Vím,~co se bude ^{E}dít, ^{E\z}býk už se v ohradě ^{Ami\z}vzpíná,

	^{Ami\z}kdo~chce, ten může ^{E\z}jít, ^{E\z}já~si dám sklenici ^{Ami\z}vína.
\end{minipage}
\begin{minipage}{0.1\textwidth}
\includegraphics[width=3cm]{../Akordy/am}
\end{minipage}

\phantom{.}

\begin{minipage}{0.65\textwidth}
\refren[2]
	^{Dmi\z}Žízeň je veliká, ^{Ami\z}život mi utíká,

	^{E\z}nechte mě příjemně ^{Ami\z}snít,

	^{Dmi\z}ve~stínu pod fíky ^{Ami\z}poslouchat slavíky,

	^{E\z}zpívat si s nima a ^{Ami\z}pít.\:\:\:\:\:
\end{minipage}
\begin{minipage}{0.1\textwidth}
\includegraphics[width=3cm]{../Akordy/e}
\end{minipage}

\phantom{.}

\begin{minipage}{0.65\textwidth}
\sloka
	^{Ami\z}Ženy jsou krásný a ^{E\z}cudný, ^{E\z}mnohá se ve mně ^{Ami\z}zhlídla,

	^{Ami\z}oči~jako dvě ^{E\z}studny, ^{E\z}vlasy jak havraní ^{Ami\z}křídla.

	^{Ami\z}Dobře vím, co znamená ^{E}pád ^{E}do nástrah dívčího ^{Ami\z}klína,

	^{Ami\z}někdo má pletky ^{E\:}rád, ^{E}já si dám sklenici ^{Ami\z}vína.
\end{minipage}
\begin{minipage}{0.1\textwidth}
\includegraphics[width=3cm]{../Akordy/dm}
\end{minipage}

\phantom{.}

\refren[2]

\sloka
	^{Ami\z}Nebe je modrý a ^{E\z}zlatý, ^{E\z}ženy krásný a ^{Ami\z}cudný,

	^{Ami\z}mantily, sváteční ^{E\z}šaty, ^{E\z}oči jako dvě ^{Ami\z}studny.

	^{Ami\z}Zmoudřel jsem stranou od ^{E\z}lidí, ^{E\z}jsem jak ta zahrada ^{Ami\z}stinná,

	^{Ami\z}kdo~chce, ať mi ^{E\z}závidí, ^{E}já si dám sklenici ^{Ami\z}vína.\:

\refren[2]

\refren[1]

\end{centerjustified}
\setcounter{Slokočet}{0}
\end{song}
\begin{figure}[h]
\predtitle\centering
\end{figure}

\end{document}
